% vim: textwidth=78 ai spell spelllang=en filetype=tex sw=2

\chapter{Stanse}\label{sec:stanse}
At the beginning of the PhD studies, we were surveying the current known
techniques used for software checking. Of course many techniques were available
and some software projects were in the development. One of the papers to
really stir the interest was about metacompilation~\cite{metacomp}. It showed
that static analysis can efficiently find many interesting errors in
real-world code. This work eventually led to a successful commercial tool
called \textsc{Coverity Prevent}~\cite{coverity}. Over the years, several
other successful tools, like \textsc{CodeSonar}~\cite{codesonar} or
\textsc{Klocwork}~\cite{klocwork} appeared.  However, such fully-featured
tools are neither free to obtain, nor is their code available (e.g.\ for
developing new algorithms or tailoring the existing tools to specific tasks).
The existing free tools were usually severely limited in what they can do,
especially when it comes to checking real-world C code such as the Linux
kernel. These include \textsc{Uno}~\cite{uno}, \textsc{Sparse}~\cite{sparse},
and \textsc{Smatch}~\cite{smatch}.

We decided to implement an open source robust framework for two purposes:
finding errors in software and to make easier development of diverse static
analysis algorithms. We named the framework \stan. It is an abbreviation of
static analyser and was deliberately chosen not to be any English word for
easier searching on the web. To fulfill \stan's role, we decided to write it in
Java as it is relatively easy to learn and is not prone to some pointer errors
class present in other languages. We tried to design it to be fully generic to
allow for implementation of many types of analyses. An implementation of such
an algorithm within the framework is called a \emph{checker}. We implemented
some checkers too, so \stan already contains four of them. The most concise
checker was inspired by the success of \textsc{Xgcc}~\cite{metacomp} and its
property-description language \textsc{Metal}~\cite{metal}. This work
implements the abstract interpretation as is discussed in
Chapter~\ref{sec:ai}.

\stan is specifically developed to check software written in C. Its parser is
powerful enough to handle the code like the whole Linux kernel. So an
implemented checker can be immediately evaluated on large real-world software.
\stan has also \emph{graphical user interface} (GUI) for easier .
This chapter covers the development of \stan including the theory behind and
report of successes and failures. \stan was released and a paper about the
framework published~\cite{stanse} some time ago. Since then, other similarly
powerful frameworks appeared.  We provide
\textsc{CPAchecker}~\cite{cpachecker} as an example.

\paragraph{Chapter Organization}
This chapter is structured as follows. Section~\ref{sec:framework} describes
the functionality provided by the framework, while Section~\ref{sec:checkers}
is devoted to the four existing checkers.  In Section~\ref{sec:results} we
present some results of running \stan on the Linux kernel. The last section
summarises the basic strengths of \stan and mentions some directions of future
development.

%}}}
%{{{ Framework

\section{Framework Functionality}
\label{sec:framework}

The \stan framework is modular and fully open.  Architecture of the framework
is depicted in Figure~\ref{fig:architecture}, and is more or less standard.
In this section we focus on the functionality of the framework itself,
describing only the non-standard or for some other reason interesting
features.
\begin{figure}[htb]
\begin{center}
    \pgfdeclarelayer{background}
    \pgfsetlayers{background,main}
    \tikzstyle{tl_uses}=[<->, dashed]
    \begin{tikzpicture}[<-, >=stealth', node distance=7mm, font=\sffamily]
      \node [tl_state] (AST) {AST};
      \node [tl_state, right=2mm of AST] (CFG) {CFG};
      \node [tl_state, right=2mm of CFG, align=center] (CG) {Call \\ Graph};
      \node [tl_state, right=2mm of CG] (MORE) {\dots};
      \begin{pgfonlayer}{background}
	\node [tl_state, inner xsep=3mm, inner ysep=5mm, yshift=3mm,
	  bottom color=black!10, top color=black!2,
	  fit= (AST) (CFG) (CG) (MORE)] (INT) {};
	\node [below=1.5mm of INT.north] {Internal Representation};
      \end{pgfonlayer}

      \node [tl_phase, below=4mm of INT] (PARSER) {Parser}
        edge[->] node {} (INT);
      \node [tl_state, align=center, below=4mm of PARSER] (PREPF) {
	  Preprocessed \\
	  Files
	}
        edge[->] node {} (PARSER);
      \node [tl_phase, left=of PREPF] (PREP) {Preprocessor}
        edge[->] node {} (PREPF);
      \node [tl_state, left=of PREP] (SRC) {Sources}
        edge[->] node {} (PREP);

      \node [tl_phase, align=center, left=of INT] (NAVIG) {
	  Code \\
	  Navigation \\
	  Libraries
	}
	edge[tl_uses] node {} (INT);
      \node [tl_phase, align=center, text width=8cm, above=of INT, xshift=-1cm]
	(INTF) {Checker Interface}
	edge[<->, in=90, out=250] node {} (NAVIG)
	edge[in=90, out=290] node {} (INT);
      \node [tl_phase, above=of INTF, xshift=-2cm] (CH1) {Checker 1}
	edge[<->, in=110, out=270] node {} (INTF);
      \node [tl_phase, above=of INTF, xshift=2cm] (CH2) {Checker N}
	edge[<->, in=70, out=270] node {} (INTF);
      \draw[-, loosely dotted] (CH1) -- (CH2);
      \node [tl_state, align=center, right=of CH2] (CH2D) {Checker N\\ Data File}
	edge[->] node {} (CH2);

      \node [tl_state, align=center, right=of INTF] (REP) {Error \\ Reports}
	edge node {} (INTF);
      \node [tl_phase, align=center, below=of REP] (GUI) {
	  GUI \\
	  Error \\
	  Report \\
	  Tracing
	}
	edge node {} (REP)
	edge[<->, in=0, out=270] node {} (PREPF);

    \end{tikzpicture}
\end{center}
  \caption{\stan framework architecture.}
  \label{fig:architecture}
\end{figure}

\subsection*{Configuration}

The \stan framework contains data structures capturing a configuration of an
analysis to be performed. We always need to know what source files to analyse
and by what checkers. That information we call \emph{configuration}. 

There are several ways how to tell \stan which source files to analyse. Besides
the standard possibilities (a given file, all files from a given
directory, all files listed in a list), \stan can also derive all
the necessary information from a project \code{Makefile}. In this case,
\stan also remembers compiler flags for preprocessing purposes. This
functionality is inspired by the \textsc{Sparse} tool~\cite{sparse}.

The user must also specify checkers which should be run on the chosen source
files.  There can be many checkers running simultaneously in \stan.
Nonetheless when they want to share data, thy cannot proceed independently.
Checkers themselves can also be configured through their own configuration
files as can be see in Figure~\ref{fig:architecture} for \textsf{Checker N}.
The configuration can be passed to the \stan framework either via command line
arguments, or using the graphical interface.

\subsection*{Parsing Source Files}

The \stan framework can process source files written in C, more
specifically in the ISO/ANSI C99 standard~\cite{ansi_c} together with most of
the GNU C extensions~\cite{gcc}. This allows \stan to process software
projects like the Linux kernel. We also provide an experimental implementation
to support C++~\cite{ansi_cpp}.  The parsing pipeline includes preprocessing
of source files as can be seen at the bottom of Figure~\ref{fig:architecture}.
This is done by the standard GNU C preprocessor \textsc{Cpp}.  It is important
to note that we do not parse the source files in a sequence one by one as they
appear in a configuration. Since we use streaming (which we discuss later),
the pipeline is applied as on-demand for each individual source file.

The parser used in \stan is generated using the \textsc{AntLR}
tool~\cite{antlr4} from our own annotated C grammar. The reason for us to
write our own parser was that at the time we started developing
\textsc{Stanse} we could not find a free parser suitable for our purposes.
The main requirement was to parse most of the GNU C as the Linux kernel makes
heavy use of GNU C language extensions. Even though there were full compilers
like \textsc{Clang}~\cite{clang}, they are not still capable to handle the
kernel fully now. And in its current form, it still contains bugs so cannot
be used reliably. Our plan is to switch from our parser to \textsc{Clang} once
it becomes stable and feature complete, or any new compiler if there appears
one.

Of course there must be at least a parser which the is able to handle the
kernel, otherwise kernel developers could not run the kernel. They use
\textsc{Gcc}.  Unfortunately parser in this compiler, and most importantly its
internal structures, is not suitable for our purposes. One reason is that
control flow graph is built on the top of register transfer language (RTL) or
tree representation. Those are very hard to work with when it comes to the
code analysis. We encourage the reader to look into the GNU C compiler
internal manual~\cite{gcc_intern} where these are described.

\subsection*{Program Internal Representation}

Once the code is parsed, it is represented using \stan's \emph{internal structures}: a call graph
among functions, a control flow graph (CFG) for each function, and an abstract
syntax tree (AST) for the whole file. Subtrees in the AST are referenced from
appropriate CFG nodes. These structures are depicted in the middle of
Figure~\ref{fig:architecture} as ``Internal Representation''. For developers
of new checkers, all these structures can be dumped in a textual or graphical
form.

Since we aim at large software projects consisting of hundreds or thousands of
modules (compilation units), it is often impossible in practice to store the
corresponding internal structures in the memory all at the same time. The \stan
framework therefore applies automatic \emph{streaming} of the internal
structures. This is currently performed on a module basis. Instead of parsing
all source modules in the beginning, a module is streamed in only when \stan
needs to access some internal structure belonging to the module. In other
words, the internal structures are constructed on demand, in a lazy manner.

If the memory occupied by internal structures exceeds a given limit, some
internal structures have to be freed before another module is streamed in. The
structures to be freed are selected using the least recently used (LRU)
approach: \stan discards all internal structures of the module whose
structures are not accessed for the longest time. If the discarded structure
is accessed again later, the corresponding module is streamed back in. Both
laziness of internal structures and streaming are completely invisible to
checkers. 

In the current implementation of streaming, each source file streamed out from
the memory is completely discarded. \stan does not back up already parsed
internal structures into auxiliary files before discarding. As a consequence,
when internal structures of the discarded file are needed again, \stan starts
the parsing pipeline of the file from scratch to recreate requested internal
structures. Although loading of previously parsed internal structures from
auxiliary files would speed up the streaming process, profiling of \stan's
performance on the Linux kernel has not shown streaming to be a performance
bottleneck. This could be easily changed if streaming performance becomes a
problem in the future, nonetheless.


\subsection*{Pointer Analysis in \stan}

Since each bug-finding/program analysis technique may have different
requirements regarding pointer analysis, the \stan framework currently
provides a developer with two pointer analyses:
\textit{Steensgaard}'s~\cite{steensgaard} and
\textit{Shapiro-Horowitz}'s~\cite{shapiro_horwitz}. See
Section~\ref{sec:pointer_analysis} for introduction to pointer analyses or
Hind's discussion~\cite{ptr_solved}.

\subsection*{Matching Language Constructs}

Many static analyses change their internal state only on some subset of
program expressions. For example, when finding race conditions in a
parallel program, one may only focus on expressions involving synchronisation,
while ignoring all others. The \stan framework therefore provides a
specification language for determining a set of program expressions.

The language defines a collection of \emph{patterns}. Each pattern is supposed
to identify a single specific kind of sub-trees in an AST of an analysed
program.  A pattern itself is therefore also a sub-tree of the AST, where some
of its vertices are ``special''. The special vertices allow to define a set of
possible sub-trees at that vertex, e.g.\ wildcards.

This is, nonetheless, not the only possible approach. For example, in the
METAL~\cite{metal} specification language, a C expression can be directly
parametrised to define a set. The solution we implemented exploits the fact
that checkers in \stan work with the AST intensively, and therefore
identifying expressions in terms of the AST is very practical.

\subsection*{Traversing Internal Representation}\label{sec:travers}

Although a checker may need to work with the internal structures in an
arbitrary way, most checkers walk through CFGs using some standard
strategy. To prevent unnecessary reimplementations, the most important and
heavily used traversal methods are implemented directly inside the framework.
With this functionality, one can implement a new checker (or its part) by
specifying
\begin{itemize}
\item whether it should go through CFGs forwards or backwards, breadth-first
  or depth-first,
\item whether the walk-through should be performed inter-procedurally or not
  (if not, the function calls are ignored), and
\item a method (callback) to be called for each visited node in a CFG.
\end{itemize}
This makes implementation of new algorithms extremely simple.

For example, when a checker needs to implement a forward flow-sensitive
analysis, it may ask the \stan framework to traverse paths in CFGs in forward
depth-first manner. This can be implemented by a single call to a function
\code{traverseCFGToDepthForward}, which takes as an argument a CFG and a
subclass of \stan's class \code{CFGPathVisitor}. In this class the checker
defines the action which should be taken whenever a CFG node is visited
(already in the requested order). The checker implements the action in a
method \code{visit} of the subclass.

In addition, for those inter-procedural analyses which do not construct
summaries, \stan provides an automated traversal among different CFGs
according to function calls (involving automated parameters passing and
value returning). Again, this can be done using a single call to the \stan
framework.  

All this functionality described in this section is shown in
Figure~\ref{fig:architecture} as ``Code Navigating Libraries''.

\subsection*{Support for Function Summaries}\label{ssec:summaries}

Inter-procedural analyses typically build function summaries. Unfortunately,
these summaries may differ from one analysis to another. Nonetheless, quite
common part in building many summaries is passing formal and actual parameters
to call sites and mapping return values to appropriate variables. Therefore,
the \stan framework provides classes simplifying the parameter passing and
values returning for checkers. These classes also provide a conversion of a
given expression in the caller into an equal expression in the called
function. The conversion can also be requested in the opposite direction,
i.e.\ for returned values.

\subsection*{The Concept of Checkers}

In the \stan framework, a checker is an implementation of some concrete
static analysis technique. Each checker has an access to shared internal
structures of analysed source files. They are also provided with an access to
the algorithms providing navigation in those structures. This is done by an
interface among checkers, internal structures and libraries of the framework.
The interface is depicted in Figure~\ref{fig:architecture} right bellow the
checkers.

The checkers are integrated in the framework of \stan using concrete factory
design pattern. Therefore, to insert a new checker to the framework one needs
to implement generic checker interface and register it to the checkers'
factory of the framework. Then it gains a full access to the features of the
framework accessible through the discussed interface.  The process of
integration requires only three simple steps. The first step is to create a
subclass of \stan's abstract class \code{Checker}, say \code{MyChecker}.
The most important method to implement is \code{check}. This is a method where
the analysis algorithm should be implemented. The second step is to integrate
the newly created class into the framework. This means implementation of
\code{MyCheckerCreator}, a subclass of the \code{CheckerCreator} abstract
class. The final step is to register the class \code{MyCheckerCreator}. It
comprises adding a line \code{registerCheckerCreator(new MyCheckerCreator())}
at the end of \code{CheckerFactory.java}.

\subsection*{Processing Errors}

\begin{figure}[t]
  \centering
    \includegraphics[width=12cm]{images/screen-bug-tracing}
  \caption{Error trace browser in the \stan GUI.}
  \label{fig:bugtrace}
\end{figure}

Once a checker finds an error, it reports the error back to the framework in
the form of an annotated \emph{error trace} -- a path in the analysed code
demonstrating this error. We provide a datatype in the \stan framework to
describe an error. The framework yields three possibilities how to present
the error traces back to the user of \stan. First, they can be printed to the
console. But as it is only a textual output, the traces are not very usable
for large projects. Hence the GUI in \stan has a built-in trace browser which
supports displaying the traces.  There is an example of a trace in the GUI in
Figure~\ref{fig:bugtrace}. Finally, the traces can be saved to an external
file in the XML format. This XML file has a wide variety of possible
applications. For example, we supply a tool transforming the XML file into an
\textsc{SQLite} database. The database is supplemented with a web interface
allowing to browse errors in the database via a web browser. Using the web
browser or the built-in graphical error browser, one can mark errors as real
errors or false positives. \stan also provides various statistics of errors
like number of errors per checker, frequency of errors of the same kind,
percentage of false positives (based on the user feedback).

Error reporting and tracing pipeline is depicted to the right of interface and
internal representation in Figure~\ref{fig:architecture}.

%}}}
%{{{ Checkers

\section{Checkers}\label{sec:checkers}

In this section we describe the four currently available checkers. The first
three checkers need a configuration of project-specific parameters like
an enumeration of special function calls.  We provide a sample configuration for
all of them so that they can be used out of the box. Of course, the
configuration can be tailored to specific needs of other projects easily.

\subsection*{\AC}
The \AC is heavily influenced by \textsc{Xgcc}~\cite{metacomp}. As can be seen
from the name of the checker, it employs finite state automata. They describe
the properties we want to check. These include checking locking discipline,
interrupt management, null pointer dereference, dangling pointers and many
others. Given such automaton, the checker tries to match the automata against
code execution. To perform that, the \AC needs to know when to perform a
transition. This is accomplished by specifying code patterns on the edges.
Since the checker works on the control flow graph and we have subtrees of
abstract syntax trees in there, patterns are AST subtrees too. It works
similar to pattern matching described in Section~\ref{sec:pattern_matching}:
we have an AST subtree with possible wildcards. An example of an automaton
for checking proper pairing of locks is presented in
Figure~\ref{fig:automaton}.
\begin{figure}[htb]
  \begin{center}
    \tikzstyle{stat} = [circle,thick,draw,minimum size=1cm]
    \tikzstyle{edge} = [shorten <=1pt,>=stealth',semithick]
    \begin{tikzpicture}[node distance=2cm,font=\sffamily]
        \node [stat] (ERR) {\textbf{Err}}
	  edge [edge,loop above] node[right, xshift=2mm, yshift=-2mm]
	  {\code{lock()}/\code{unlock()}} ();
        \node [stat] (U) [below left=of ERR] {U}
	  edge [edge,->] node[above left] {\code{unlock()}} (ERR);
        \node [stat] (L) [below right=of ERR] {L}
	  edge [edge,->] node[above right] {\code{lock()}} (ERR)
	  edge [edge,<-,bend right=20] node[above] {\code{lock()}} (U)
	  edge [edge,->,bend left=20] node[below] {\code{unlock()}} (U);
	\node [left=6mm of U] {}
	  edge [edge,->] node {} (U);
    \end{tikzpicture}
  \end{center}
  \caption{An automaton for checking locking properties.}
  \label{fig:automaton}
\end{figure}
The automaton starts in the unlocked state (\textsf{U}). A transition from
that state is fired either when the checker matches a call to \code{unlock} or
\code{lock}. In the former case, the automaton is transferred to the error
state (\textsf{Err}) and the checker flags an error by generating an error
trace. In the other case, the checker simply transfers the automaton to the
locked state (\textsf{L}). The evaluation continues until a fixed point is
reached.

Compared to the implementation described by Engler et al.~\cite{metacomp,
metal}, our technique differs in several aspects. In particular, we do not use
metacompilation, automata are not input-language specific (a pattern matching
is used instead), and the inter-procedural analysis is done in the context of
a single input file.

\subsection*{\LC}
The \LC tries to find locking errors in analyzed code. It performs so by
accounting statistics about locked accesses to a variable. This approach is
not new and is based on a previous work~\cite{lockchecker}.  Apart of tracking
accesses to a variable, it also tracks which locks are locked while each
of the variable is accessed. Then the \LC combines the two pieces of
information together and counts statistics in how many cases each variable is
accessed while some lock is held. If the difference is proportional, an error
is reported. For instance if, some variable is changed 99 times while some
lock is held and in one case the lock is not held, this is reported as a
possible error. The boundary is currently set to 70. So if there is a variable
accessed under a lock in more than 70\% cases and once or more it is not, the
latter cases are reported as errant. Similar to the \AC, both variable
accesses and locks are specified in a configuration by patterns.

\subsection*{\TC}
The \TC is aimed at checking of possible deadlocks in concurrent programs.
The technique is based on the notions of locksets~\cite{relay} and deadlock
detection by looking for cycles in \emph{resource allocation graphs} (RAG).
The \TC first tries to identify the parts of code which can run in parallel,
as different threads. This is performed by searching of functions
instantiating threads. We implicitly try functions such as
\code{pthread\_create}. This is not enough for projects like the Linux kernel.
There need not be explicit thread creation functions in modules. As some hooks
in modules can be called in parallel, the code can run concurrently
nonetheless. So on the top of list of thread creation functions, a user can
also specify manually which hooks can run concurrently.

With the knowledge of where a thread starts, the \TC builds a set of
\emph{dependency graphs} for each such thread. A dependency graph statically
represents possible locksets during one execution of a thread. Dependency
graphs are then combined and transformed into RAGs. If there is a circular
lock dependency, a RAG contains a cycle. In such case an error is reported to
the \stan framework.

\subsection*{\RC}
The \RC is the simplest and also the smallest checker in \stan. Its name
correctly suggests searching CFGs for unreachable nodes. These are then
reported as warnings or errors depending on importance. The importance stems
from the nature of an unreachable node. For example superfluous semicolons
(which are still statements) are less important than an unused branch. The
primary goal of the \RC is to demonstrate the simplicity of a new checker
implementation. With help of the framework features described in
Section~\ref{sec:travers}, the code of the checker has less than 200 lines
including the mentioned error/warning classification and comments.  Even
though it is a very simple checker, it was still able to find serious bugs in
the Linux kernel. For example a superfluous semicolon can cause unexpected
unconditional returns from functions like in the code of
Figure~\ref{fig:stan_unreach}. As we can see on line~\ref{line:stan_unreach},
there is a superfluous semicolon. It causes a preliminary unconditional return
from a function. This was an error in the \code{ath9k} wireless driver.
\begin{figure}[htb]
\begin{center}
\begin{tabular}{c}
\begin{lstlisting}
if (rate_table->info[rate].valid_single_stream &&
    !(ath_rc_priv->ht_cap & WLAN_RC_DS_FLAG));(*@\label{line:stan_unreach}@*)
  return rate;
\end{lstlisting}
\end{tabular}
\end{center}
  \caption{A bug found by the \RC.}
  \label{fig:stan_unreach}
\end{figure}

%}}}
%{{{ Results

\section{Results on the Linux Kernel}\label{sec:results}

We have several reasons to choose the Linux kernel for testing \stan: the
kernel is a large and freely available codebase, it fully exercises most
of the features of ISO/ANSI C99~\cite{ansi_c} and GNU C extensions~\cite{gcc},
it is under constant development (there is a constant income of new bugs), and
the absence of bugs is of a great concern.  We applied \stan, together with
the four checkers described in Section~\ref{sec:checkers}, to the Linux kernel
version 2.6.28. The \AC was configured with four automata describing the
following types of errors:
\begin{itemize}
\item incorrect pairing of functions (imbalanced locking, reference
  counting errors),
\item resource or memory leaks,
\item errors in pointer manipulation (null dereference, dangling pointers,
  etc.), and 
\item deadlocks caused by sleeping inside spinlocks or interrupt handlers.
\end{itemize}

The running time of \stan on a common desktop machine with two 2.5\,GHz
cores and 4\,GiB of memory was under two hours. The memory usage of the Java
process oscillated between 400 and 1000\,MiB. The number of reports by the
checkers is presented in Table~\ref{tab:results_stan}.
\begin{table}[tb]
\begin{center}
%\setlength{\tabcolsep}{3pt}
\begin{tabular}{|ll||c|c|c|c|} \hline
\multicolumn{2}{|c||}{\multirow{2}{*}{\textbf{Checker}}} &
  \multicolumn{3}{c|}{\textbf{Reports}} & \textbf{Real/classified} \\ \cline{3-5}
&&\textbf{Total} & \textbf{Real} & \textbf{False} & \textbf{Error Ratio} \\ \hline\hline
\multirow{4}{6em}{\textsc{Automaton} \textsc{Checker}}
& \multicolumn{1}{|l||}{Pairing}   & 543 & 73  & 441 & 14.2\,\% \\ \cline{2-6}
& \multicolumn{1}{|l||}{Leaks}     & 108 & 20  & 86  & 18.9\,\% \\ \cline{2-6}
& \multicolumn{1}{|l||}{Pointers}  & 35  & 17  & 18  & 48.6\,\% \\ \cline{2-6}
& \multicolumn{1}{|l||}{Deadlocks} & 35  & 16  & 19  & 45.7\,\% \\ \hline
\multicolumn{2}{|l||}{\LC}         & 13  & 6   & 7   & 46.2\,\% \\ \hline
\multicolumn{2}{|l||}{\TC}         & 20  & 9   & 11  & 45.0\,\% \\ \hline
\multicolumn{2}{|l||}{\RC}         & 31  & 31  & 0   & 100.0\,\% \\ \hline\hline
\multicolumn{2}{|l||}{\textbf{Overall}} & \textbf{785} & \textbf{172} &
  \textbf{582} & \textbf{22.8\,\%} \\ \hline
\end{tabular}
\end{center}
\caption{\stan results on the Linux kernel version 2.6.28.}
\label{tab:results_stan}
\end{table}
Let us note that \RC actually found 751 errors, but 720 of them are of low
importance (including 696 superfluous semicolons). Even though they are real,
we omitted them from the statistics.

We have manually analysed the reports and tried to classify them as real
errors or false positives. We classified all except 31 reports found by \AC.
There, we are not able to decide whether it is a real error or not in a short
time. The numbers of real errors, false positives and the ratio of real errors
to all classified errors can be also found in Table~\ref{tab:results_stan}.
The overall ratio of real errors to all classified errors is not high: 22.8\%.
Nonetheless, \stan in the current version does not have any thorough false
positive filtering technique, which may be implemented in the future. Finally,
note that the checkers do not produce any false negatives assuming there is no
bug in the checkers' implementation. The reason is that all the checkers
implement may analyses, over-approximating the set of error behaviours.

More than 120 of the 172 real errors were reported to kernel developers and
fixed in the following kernel releases. The rest was either discovered by
independent checkers and reported by someone else or the incorrect code
disappeared from the kernel completely before we finished our evaluation of
the reports. Interestingly, some of the reported errors remained undiscovered
for more than seven years. For illustration, see our report at
\url{http://lkml.org/lkml/2009/3/11/380}. The whole list of reports generated
by \stan on that particular kernel version can be found in our database from
Chapter~\ref{sec:clabure}.

\subsection*{Important Errors Found by \stan}
Although checkers currently implemented in \stan are based on widely known
techniques, running them on the Linux kernel helped to uncover several
important bugs. In the text below we present two typical bugs discovered by
\stan, each using a different checker.

\paragraph{AutomatonChecker}
Many bugs found by the \AC trigger only under specific conditions, however
some of them may be visible to the user. Consider this code excerpt taken
from the 2.6.27 kernel, \textsf{drivers/pci/hotplug/pciehp\_core.c} file,
\textsf{set\_lock\_status} function:

\begin{lstlisting}[language=C]
mutex_lock(&slot->ctrl->crit_sect);
/* has it been >1 sec since our last toggle? */
if ((get_seconds() - slot->last_emi_toggle) < 1)
	return -EINVAL;
\end{lstlisting}

Note that the call to \textsf{mutex\_lock} function is followed by an \textbf{if}
statement, which returns immediately in the true branch, omitting a call to 
\textsf{mutex\_unlock}. In fact this
deadlock could be easily triggered by a user. It is sufficient to write "1" to
\textsf{/sys/bus/pci/slots/.../lock} file twice within a second.

\paragraph{ThreadChecker}
%http://lkml.org/lkml/2009/4/14/527
An example of non-trivial error which could not be found by the \AC. The code
described here is from the 2.6.28 kernel, file \textsf{fs/ecryptfs/messaging.c}.

There are three locks in the code,
\textsf{msg\_ctx-\textgreater mux}, which is local per context, and two
global locks --
\textsf{ecryptfs\_daemon\_hash\_mux} and
\textsf{ecryptfs\_msg\_ctx\_lists\_mux}.

Let us denote lock dependencies as a binary relation where the first component
depends on the second. I.e. \textsf{lock(A)} followed by \textsf{lock(B)}
means dependency B on A, and we write $A\leftarrow B$.

\begin{lstlisting}[language=C,numbers=left,escapeinside={XY}{YX}]
int ecryptfs_process_response(...)
{
	...
	mutex_lock(&msg_ctx->mux); XY\label{ecrypt_mux1}YX
	mutex_lock(&ecryptfs_daemon_hash_mux); XY\label{ecrypt_hmux1}YX
	...
	mutex_unlock(&ecryptfs_daemon_hash_mux);
	...
unlock:
	mutex_unlock(&msg_ctx->mux);
out:
	return rc;
}
\end{lstlisting}
Here the two locks on lines~\ref{ecrypt_mux1} and~\ref{ecrypt_hmux1} give
\textsf{msg\_ctx-\textgreater mux} $\leftarrow$
\textsf{ecryptfs\_daemon\_hash\_mux}.
\begin{lstlisting}[language=C,numbers=left,escapeinside={XY}{YX},firstnumber=last]
static int ecryptfs_send_message_locked(...)
{
	...
	mutex_lock(&ecryptfs_msg_ctx_lists_mux); XY\label{ecrypt_lmux2}YX
	...
	mutex_unlock(&ecryptfs_msg_ctx_lists_mux);
	...
}

int ecryptfs_send_message(...)
{
	int rc;

	mutex_lock(&ecryptfs_daemon_hash_mux); XY\label{ecrypt_hmux2}YX
	rc = ecryptfs_send_message_locked(...) XY\label{ecrypt_call}YX
	mutex_unlock(&ecryptfs_daemon_hash_mux);
	return rc;
}
\end{lstlisting}
At line~\ref{ecrypt_call}, function \textsf{ecryptfs\_send\_message\_locked}
is called from \textsf{ecryptfs\_send\_message}, hence the locks at
lines~\ref{ecrypt_lmux2} and~\ref{ecrypt_hmux2} generate lock dependency of
\textsf{ecryptfs\_daemon\_hash\_mux} $\leftarrow$
\textsf{ecryptfs\_msg\_ctx\_lists\_mux}.
\begin{lstlisting}[language=C,numbers=left,escapeinside={XY}{YX},firstnumber=last]
int ecryptfs_wait_for_response(...)
{
	...
	mutex_lock(&ecryptfs_msg_ctx_lists_mux);
	mutex_lock(&msg_ctx->mux);
	...
	mutex_unlock(&msg_ctx->mux);
	mutex_unlock(&ecryptfs_msg_ctx_lists_mux);
	return rc;
}
\end{lstlisting}
Finally, this function introduces \textsf{ecryptfs\_msg\_ctx\_lists\_mux}
$\leftarrow$ \textsf{msg\_ctx-\textgreater mux}.

Composing these results together the following circular dependency of these
three locks was found:
\begin{itemize}
\item \textsf{msg\_ctx-\textgreater mux} $\leftarrow$ \textsf{ecryptfs\_daemon\_hash\_mux}
\item \textsf{ecryptfs\_daemon\_hash\_mux} $\leftarrow$ \textsf{ecryptfs\_msg\_ctx\_lists\_mux}
\item \textsf{ecryptfs\_msg\_ctx\_lists\_mux} $\leftarrow$
  \textsf{msg\_ctx-\textgreater mux}
\end{itemize}

This issue was later confirmed as a real bug leading to a
deadlock\footnote{http://lkml.org/lkml/2009/4/14/527}.

%}}}
%{{{ Conclusions

\section{Conclusions and Future Work}

\stan is a free Java-based framework design for simple and efficient
implementation of bug-finding algorithms based on static analysis. The
framework can process large-scale software projects written in ISO/ANSI C99,
together with most the GNU C extensions. \stan does not currently use any new
techniques -- its novelty comes from the fact that (to our best knowledge)
there is no other open-source framework with comparable applicability and
efficiency. We note that more than 130 bugs found by \stan have been reported
to and confirmed by Linux kernel developers already. More information and the
tool itself can be found at \url{http://stanse.fi.muni.cz/}.

\subsubsection{Future Work}
We plan to improve the framework in several directions. Firstly we are
currently working on C++ support. Furthermore we plan to provide \stan in the
form of an IDE plug-in, e.g.~for Eclipse and NetBeans. A lot of work can be
done in the area of automatic false alarm filtering and error importance
classification.  Independently of developing new features, we would like to
speed up the framework as well. To this end we intend to replace the current
parser written in Java by an optimised parser written in C, to replace the XML
format of internal structures by a more succinct representation, to add a
support for function summaries, etc.

\section{Results}

Example of false-negatives/positives -- where it does not work, description of
where it does not work.
